% Part: computability
% Chapter: computability-theory
% Section: non-comp-set

\documentclass[../../../include/open-logic-section]{subfiles}

\begin{document}

\olfileid[fr]{cmp}{thy}{ncp}
\olsection{Il existe des ensembles non calculables}

Nous avons vu que tout ensemble calculable est
effectivement énumérable. La réciproque est-elle
vraie ? Le résultat suivant montre que non.

\begin{thm}\ollabel{thm:K-0}
Soit $K_0$ l'ensemble
$\Setabs{\tuple{e, x}}{\cfind{e}(x) \fdefined}$.
Alors $K_0$ est effectivement énumérable,
mais non calculable.
\end{thm}

\begin{proof}
Pour voir que $K_0$ est effectivement énumérable,
remarquons qu'il est le domaine de la fonction~$f$
définie par
\[
f(z) = \umin{y}{(\fn{SeqCode}(z) \land \len{z} = 2 \land T((z)_0, (z)_1, y))}.
\]
En effet, si $\cfind{e}(x)$ est définie,
$f(\tuple{e,x})$ trouve le code d'un calcul qui
s'arrête. Si $\cfind{e}(x)$ n'est pas définie,
$f(\tuple{e,x})$ ne l'est pas non plus. Enfin,
si $z$ ne code même pas une paire, $f(z)$ est
également indéfinie.\footnote{Correction éditoriale : nous ajoutons le test $\fn{SeqCode}(z)$ défini au chapitre précédent. La longueur calculée seule n'exclut pas un entier non codant : multiplier le code d'une paire par $7$, sans facteur $5$, conserve cette longueur et les deux coordonnées, mais ne donne plus un code de suite.}

La non-calculabilité de $K_0$ n'est autre que
l'indécidabilité du problème de l'arrêt,
\olref[hlt]{thm:halting-problem}.
\end{proof}

L'ensemble $K_0$ rassemble les paires
$\tuple{e,x}$ telles que $\cfind{e}(x)\fdefined$;
autrement dit, $\tuple{e,x}\in K_0$ si et
seulement si $\cfind{e}$ s'arrête sur l'entrée~$x$.
On l'appelle aussi l'« ensemble de l'arrêt ».
L'ensemble
$K=\Setabs{e}{\cfind{e}(e)\fdefined}$ est
l'« ensemble de l'arrêt diagonal ». Il sert souvent
d'exemple canonique d'ensemble indécidable.

\begin{thm}\ollabel{thm:K}
L'ensemble de l'arrêt diagonal
$K=\Setabs{e}{\cfind{e}(e)\fdefined}$ est
!!{c.e.}, mais non décidable.
\end{thm}

\begin{proof}
Supposons $K$ décidable, c'est-à-dire sa fonction
caractéristique $\Char{K}$ calculable. Posons
\[d(e) = \begin{cases}
  1 & \text{si\/ $\Char{K}(e) = 0$}\\
  \fundefined & \text{sinon.}
  \end{cases}
  \]
Cette fonction est calculable partielle sous l'hypothèse faite : $d(e)\simeq 1+\umin{y}{\Char{K}(e)=0}$. Soit $k$ un indice de~$d$, donc
$d\simeq\cfind{k}$. On a alors
$d(k)\simeq\cfind{k}(k)$, contrairement au
fait que $d(k)\fdefined$ si et seulement si
$\cfind{k}(k)\fundefined$, qui découle
de la définition de~$d$.

$K$ est le domaine de
$f(x)=\umin{y}{T(x,x,y)}$;
il est donc !!{c.e.}
\end{proof}

\end{document}
